\documentclass[11pt]{article}
\usepackage{setspace}
\usepackage{geometry}
\geometry{margin=1in}
\usepackage{graphicx}
\setstretch{1.15}

\begin{document}

\begin{flushright}
\textbf{Zhongli Wang} \\
% Phone: +49 160 9541 6015 \\
% Email: nkwangzhongli@gmail.com \\
% Website: https://sites.google.com/view/zhongli-wang \\
Astastrasse 22 \\
33617 Bielefeld \\
Germany \\
\end{flushright}

% \vspace{0.45cm}

\noindent
\textbf{Prof. Dr. Irene Monasterolo} \\
Institute for Ecological Economics \\
Vienna University of Economics and Business (WU) \\
Austria \\

% \vspace{0.45cm}

\noindent
\textbf{Application for Post-Doctoral Researcher --- Climate and Nature Finance (ID 2922)}\\

\vspace{0.25cm}

\noindent Dear Professor Monasterolo,\\[0.15cm]

I am writing to apply for the post-doctoral researcher position in climate and nature finance at the Institute for Ecological Economics, WU Vienna. 
I recently defended my Ph.D.\ in Economics at Bielefeld University and Universit\`a Cattolica del Sacro Cuore (Marie Sk{\l}odowska--Curie EPOC programme). 
My research focuses on complexity in economic systems, for instance, 
how heterogeneous agents interact through networks, how incentives and information shape strategic behavior, and how institutional interventions change system-level outcomes. 
The project's aim of informing how financial arrangements and instruments can build resilience to climate and nature risks for EMDEs aligns directly with my interest in building empirically grounded, 
complexity-informed models that can explain and improve real-world policy outcomes.

My doctoral work develops agent-based models of market environments where outcomes are non-linear, path-dependent, and shaped by feedback loops. 
In my dissertation, \emph{Strategic Behavior and Economic Dynamics in ICT Markets: An Agent-Based Framework}, 
I study how firms and heterogeneous users interact under network effects, congestion constraints, and imperfect information. 
Methodologically, this work required (i) formalizing behavioral and strategic decision rules, 
(ii) building scalable simulation pipelines in Python, and (iii) conducting counterfactual analysis to evaluate policy-relevant interventions. 
These skills translate directly to the project's objectives: developing Stock-Flow Consistent behavioral models and stress-test models for macro-financial risk assessment of climate and nature risks requires the same combination of agent heterogeneity, 
feedback mechanisms, and policy counterfactual analysis that I have practiced throughout my doctoral research.

A central contribution I can make to the project is to connect three components into a coherent research pipeline: 
(1) constructing datasets on financial instruments and their links to climate and nature risks in EMDEs, 
(2) network identification of key channels of financial dependence and vulnerability, 
and (3) an agent-based or stock-flow consistent model that can stress-test financial instruments, 
such as the Resilience and Sustainability Facility, Capital Flow Management Measures, or a global insurance scheme, 
under realistic behavioral responses. 
I have substantial experience operationalizing network concepts in economic settings: 
defining nodes and edges from messy sources, constructing interpretable measures of influence and centrality, 
and translating network structure into behavioral mechanisms in simulations. 
This would support the identification of financial institutions and arrangements whose lending, underwriting, 
or advisory services are most central for climate resilience in EMDEs, providing a rigorous basis for policy prioritization.

I am comfortable working with large, heterogeneous datasets and building reproducible data infrastructure in Python. 
In collaborative work on \emph{Polytope Fraud Theory} (published in the \emph{International Review of Financial Analysis}, 2025), 
I used supervised machine learning to detect structured anomalies in financial statements, 
strengthening my experience in feature engineering, robustness checks, and translating algorithmic signals into economic interpretation. 
In a second project, \emph{Agentic Workflows for Economic Research} (arXiv, 2025), 
I helped develop a framework for integrating large-language-model agents into empirical pipelines for document collection and classification. 
For the project, these experiences are directly relevant for constructing and maintaining the project database, documenting extraction and coding procedures, and ensuring that data and results remain auditable and reusable for academic publications and policy outputs.

% My near-term plan would be to deliver publishable outputs through a staged strategy aligned with the project's deliverables. 
% First, I would help build the core dataset linking financial instruments to climate and nature risks in EMDEs, with careful attention to measurement validity and reproducibility (including codebooks, versioning, and transparent data provenance). 
% Second, I would implement network analyses to identify key channels of financial dependence and vulnerability (e.g., debt structures, capital flow patterns, and concentration of external financing), producing a first co-authored manuscript focused on mapping and measurement. 
% Third, I would develop a stock-flow consistent or agent-based model that combines institution heterogeneity (business models, risk constraints, and regulatory incentives) with country heterogeneity (financing needs, climate exposure, and fiscal space) to evaluate the effects of alternative financial instruments. 
% The model would be designed to produce testable implications and clear policy counterfactuals, supporting a second manuscript centered on mechanisms and instrument design.

Beyond technical fit, I am motivated by the project's interdisciplinary and policy-facing setting. 
My background includes economic training and practical experience, 
which gives me an understanding of monetary policy instruments, banking operations, and financial stability considerations relevant to the project's focus on EMDEs. 
I would actively contribute to the Task Force's dialogue with multinational development banks and financial institutions, and I value collaborative team practices with international partners. 
I am fluent in English (C1) and have working proficiency in German (C1) and Mandarin (native), which facilitates engagement with diverse stakeholders. 
I am available to travel internationally for conferences and meetings.

% The opportunity to work with the Task Force on Climate, Development and the International Financial Architecture---in dialogue with the G24, V20, and multinational development banks---is particularly appealing. I see this position as a chance to deliver high-quality publications while building a professional profile at the intersection of climate finance, macro-financial modeling, and policy-relevant research for development.

Thank you so much for considering my application! I would welcome the opportunity to discuss how my experience and abilities can contribute to the project's mission and to WU Vienna's research community.

\vspace{0.5cm}

\noindent Sincerely,\\[0.3cm]
\noindent\includegraphics[width=2.7cm,keepaspectratio]{../Wang-sig.jpg}\\[-0.15cm]
\noindent Zhongli Wang

\end{document}